% PlanetMath proposal draft for arXiv:2602.23211
% Source corpora: arXiv math.CT eprints, storage/mo-processed-gpu, storage/math-processed-gpu
\section*{Coalgebraic analysis of social hypergraphs}
\subsection*{Source}
\begin{itemize}
  \item arXiv:2602.23211, "Coalgebraic analysis of social systems" (2026-02-26)
  \item StackExchange cross-links via storage/mo-processed-gpu + math-processed-gpu
\end{itemize}
\subsection*{Synopsis}
Baldan--Bonchi--Montecchiani--Pavlovic recast role/positional analysis of social systems as universal coalgebra for functors on relational, hypergraph, and higher-order interaction data. Modalities on coalgebras capture local views of actors, and bisimulations formalise when two agent configurations encode the same "role" independent of network size.
\begin{enumerate}
  \item Define a unifying functorial semantics that includes relational structures, hypergraphs, and complex network data.
  \item Introduce coalgebraic modalities interpreting positional logics for agents and interactions.
  \item Characterise behavioural equivalences via coalgebraic bisimulations and prove invariance results for social role analysis.
  \item Provide examples linking categorical semantics to applied social network datasets.
\end{enumerate}
\subsection*{MathOverflow cues}
\begin{description}
  \item[MO: coalgebraic modal logic] References on interpreting modal logics as coalgebraic predicates.
  \item[MO: categorical social networks] Threads on functorial treatments of hypergraphs and higher interactions.
\end{description}
\subsection*{Math.SE anchors}
\begin{description}
  \item[Math.SE: bisimulations for relations] Q\&A about bisimulation of relational structures and coalgebras.
  \item[Math.SE: modelling hypergraphs categorically] Discussions of polynomial functors capturing hypergraph data.
\end{description}
\subsection*{Encyclopedia outline}
\begin{enumerate}
  \item \textbf{Universal coalgebra primer.} Summarize coalgebras, modalities, and bisimulations.
  \item \textbf{Social data functors.} Describe the functors handling relational and higher-order interaction data.
  \item \textbf{Modalities for roles.} Detail how modal formulas capture positional information.
  \item \textbf{Bisimulation invariance.} Outline the categorical proofs of invariance theorems.
  \item \textbf{Worked examples.} Translate the paper's datasets into PlanetMath-ready examples.
\end{enumerate}
\subsection*{Action items}
\begin{itemize}
  \item Build a glossary entry for coalgebraic social systems referencing universal coalgebra.
  \item Prepare diagrams of typical functors (relational, hypergraph) with their coalgebra structures.
  \item Link to entries on coalgebraic modal logic, bisimulation, and applied category theory.
\end{itemize}
